\begin{textsample}{POS Dim 7 – human – Score 4.00 – mg/mg\_ef\_matematica\_5\_p14.txt}  \label{ex:f7_pos_010}
No Ensino Fundamental – Anos Finais, a expectativa é a de que os estudantes reconheçam comprimento, área, volume e abertura de ângulo como grandezas \textbf{associadas} a figuras geométricas e que consigam resolver problemas envolvendo essas grandezas com o uso de unidades de medidas padronizadas mais usuais.
Além disso, é necessário introduzir medidas de capacidade de armazenamento de computadores como grandeza associada a demandas da sociedade moderna.
Vale ressaltar que, essa unidade temática, propicia o desenvolvimento de atitudes éticas, responsáveis e sustentáveis em relação ao consumo, através de situações de compra e venda.
A unidade temática Probabilidade e Estatística trabalha com a incerteza e o tratamento de dados.
Ela propõe a \textbf{abordagem} de conceitos, fatos e procedimentos presentes em muitas situações-problema da vida cotidiana, das ciências e da tecnologia.
Assim, todos os cidadãos precisam desenvolver habilidades para coletar, organizar, \textbf{representar}, interpretar e analisar dados em uma variedade de contextos, de maneira a fazer julgamentos bem fundamentados e tomar as decisões adequadas.
Isso \textbf{inclui} raciocinar e utilizar conceitos, representações e índices estatísticos para \textbf{descrever}, explicar e predizer fenômenos.

% matched lemmas: abordagem, associar, descrever, incluir, representar
\end{textsample}
